\section{Background}
\label{sec:background}

\subsection{Cardiac Autonomic Control and Heart Rate Variability}
\label{sec:bg-autonomic}

\subsubsection{The heart under central control}

The sinoatrial node sets the pace of the heart, but its rate is continuously adjusted by the two branches of the autonomic nervous system. Sympathetic outflow accelerates the heart and increases contractility through noradrenergic signalling with a latency of several seconds. Parasympathetic outflow through the vagus nerve slows the heart through acetylcholine acting on muscarinic receptors, with a latency short enough to influence the very next beat \citep{taskforce1996}. Because the two branches act on different time scales, the beat-to-beat timing of the heart carries information about both. Fast fluctuations are dominated by vagal activity, whereas slower fluctuations reflect a mixture of vagal, sympathetic, baroreflex and thermoregulatory influences.

This outflow is organised centrally. \citet{benarroch1993} described a \emph{central autonomic network} of reciprocally connected structures: the insular cortex, the anterior cingulate and ventromedial prefrontal cortex, the amygdala, the hypothalamus, the periaqueductal grey, the parabrachial complex and the nucleus of the solitary tract. These structures act on the medullary nuclei that generate sympathetic and vagal output. The network integrates visceral afferent information with descending signals related to emotion, attention and behaviour, and adjusts autonomic output accordingly. Its organisation has two properties that matter for the present thesis. First, it is distributed and hierarchical, so that autonomic output at any moment reflects the combined action of many loops. Second, its cortical components overlap with regions that support attention, error monitoring and self-regulation, so that cognitive and motor demands can shape autonomic output directly, in addition to the influence of metabolic need.

The ascending side of this circuitry is equally relevant. Interoceptive accounts place the insula at the centre of the brain's representation of bodily state \citep{craig2002}. Physiological signals are first represented in the posterior insula and then integrated into consciously accessible feelings in the anterior insula. On this view the heart is both regulated by and represented in the brain, and the regulation of cardiac rhythm is one component of a broader loop linking bodily state, feeling and action.

\subsubsection{Neurovisceral integration}

The neurovisceral integration model makes the link between cardiac regulation and cognition explicit \citep{thayer2000}. It proposes that prefrontal circuits for attention and self-regulation inhibit subcortical circuits for threat and defence. The strength of this inhibition is expressed in the vagal control of the heart. High resting HRV is therefore taken to index a flexible system that adjusts quickly to changing demands. Low HRV indexes a system that struggles to leave a default defensive mode. A meta-analysis of neuroimaging studies supports the anatomical part of this account, showing that HRV is associated with activity in the ventromedial prefrontal cortex and the amygdala, among other regions \citep{thayer2012}.

For cognitive neuroscience, the attraction of this framework is practical as well as theoretical. HRV can be recorded non-invasively and continuously during a task, at a temporal resolution of single heartbeats, with equipment that does not constrain movement. It therefore offers a window on how the brain regulates the body during activities, such as object manipulation, that are difficult to study with neuroimaging.

\subsubsection{Amplitude indices of HRV and their limits}

The standard indices of HRV quantify the amount of variability. In the time domain, the standard deviation of normal-to-normal intervals (SDNN) captures overall variability. The root mean square of successive differences (RMSSD) and the percentage of successive differences larger than 50~ms (pNN50) capture beat-to-beat variability, which is mainly vagal \citep{taskforce1996}. In the frequency domain, power in the high-frequency band (0.15--0.40~Hz) tracks respiratory sinus arrhythmia and vagal modulation, while low-frequency power (0.04--0.15~Hz) reflects a mixture of influences including the baroreflex \citep{taskforce1996}. Recommendations for psychophysiological research specify recording conditions, minimum durations and the control of confounds such as caffeine, recent exercise, posture and breathing \citep{laborde2017}.

However, amplitude indices have two limitations for the present question. First, they are strongly intensity-dependent, as vagal withdrawal at the onset of physical activity drives RMSSD and high-frequency power towards a floor, so that they lose resolution as soon as a task involves movement \citep{gronwald2020biomarker}. Second, the amount of variability does not distinguish a heartbeat that varies richly across time scales from one that varies in a simple, repetitive way. A strongly periodic signal, such as the large oscillations produced by slow paced breathing, yields high RMSSD and high spectral power but little complexity. Reflex tests illustrate the point from the opposite direction. For example, resting RMSSD shares only about half or less of its variance with other markers of cardiovagal function, such as the response to paced deep breathing or to standing up \citep{blalock2026}. No single amplitude index captures autonomic regulation in full.

Taken together, HRV offers a continuous window on how the brain regulates the heart. Its amplitude indices, however, lose resolution during movement and say little about how variability is structured over time. The next section therefore turns to measures of organisation.

\subsection{Fractal Physiology and Detrended Fluctuation Analysis}
\label{sec:bg-dfa}

\subsubsection{Complexity and health}

Fractal physiology addresses the second limitation by asking how variability is organised across time scales. Many physiological signals show fluctuations that are self-similar across time scales, including the heartbeat, breathing and the stride intervals of walking \citep{goldberger2002,hausdorff1997}. Their variability over a few seconds looks like their variability over minutes once it is rescaled. A healthy heart, for example, speeds up and slows down slightly with each breath, drifts over a minute as posture and attention change, and varies over hours with activity and sleep, and the pattern of change looks similar at each of these scales. In the heartbeat of healthy young adults, these fluctuations show long-range correlations close to $1/f$ behaviour, the boundary between uncorrelated noise and strongly correlated, Brownian-like drift \citep{peng1995,goldberger2002}.

The key observation of this literature is that disease and ageing erode fractal organisation, and that they do so in two directions \citep{lipsitz1992,goldberger2002}. Some conditions push the heartbeat towards uncorrelated randomness, in which each interval carries little information about the next. Others push it towards excessive regularity, in which a single dominant rhythm takes over, as in the strongly periodic heart rate oscillations of severe heart failure. In both cases the system loses the capacity to respond across many time scales at once. On this view, physiological health corresponds to adaptable organisation across time scales, and a loss of complexity can be detected as a departure from the healthy scaling regime in either direction. Short-term fractal scaling also has clinical relevance, since a reduced short-term scaling exponent predicted mortality more strongly than traditional HRV indices in patients with depressed left ventricular function after myocardial infarction \citep{huikuri2000}.

This framework connects naturally to biological robustness. \citet{kitano2004} characterised robustness as a fundamental property of complex biological systems: the capacity to keep functioning under perturbation, through feedback control, modularity, redundancy and decoupling. A heartbeat regulated by many interacting loops operating on different time scales is, on this reading, the observable signature of a robust control architecture, and fractal scaling is one way of quantifying it. Network physiology extends the same idea to exercise, treating the organism as a network of interacting systems whose coupling reorganises with physical and cognitive demands \citep{balague2020}.

\subsubsection{Computing DFA}

Detrended fluctuation analysis was introduced to quantify correlations in non-stationary time series such as heartbeat recordings \citep{peng1995}. For a series of $N$ RR intervals $RR_i$ with mean $\overline{RR}$, the series is first integrated,
\begin{equation}
y(k) = \sum_{i=1}^{k} \left( RR_i - \overline{RR} \right), \qquad k = 1, \dots, N.
\end{equation}
The integrated series is divided into non-overlapping boxes of length $n$ beats. In each box a least-squares line $y_n(k)$ is fitted and subtracted, which removes local trends. The root-mean-square fluctuation of the detrended series,
\begin{equation}
F(n) = \sqrt{\frac{1}{N} \sum_{k=1}^{N} \left[ y(k) - y_n(k) \right]^2},
\end{equation}
is computed for a range of box sizes. If the series is fractal, $F(n)$ grows as a power of $n$, $F(n) \propto n^{\alpha}$, and the scaling exponent $\alpha$ is the slope of $\log F(n)$ against $\log n$. Consider a participant resting at about 70 beats per minute. A two-minute window then contains about 140 RR intervals, which yields 35 boxes of four beats and eight boxes of sixteen beats. If the fluctuation function doubles when the box size quadruples from 4 to 16 beats, the slope in log--log coordinates is 0.5, the value expected for uncorrelated noise; if it quadruples, the slope is 1.0. Heartbeat series typically show a crossover between short and long time scales, so two exponents are estimated: $\alpha_1$ over box sizes of 4 to 16 beats and $\alpha_2$ over 16 to 64 beats \citep{peng1995}. The present thesis concerns $\alpha_1$, which can be estimated from short recordings and changes within minutes.

\subsubsection{Interpreting $\alpha_1$}

What does a given value of \al{} mean physiologically? The value of $\alpha$ maps onto classes of signal. An uncorrelated series (white noise) gives $\alpha = 0.5$, a series with $1/f$ correlations gives $\alpha = 1.0$, and Brownian motion, the integral of white noise, gives $\alpha = 1.5$. Values between 0.5 and 1.0 indicate persistent but stationary correlations, and values between 1.0 and 1.5 indicate increasingly non-stationary, drift-dominated behaviour. \citet{hoos2025} argue that these values should be read against two classes of signal: fractional Gaussian noise (fGN, $\alpha < 1$) and fractional Brownian motion (fBM, $\alpha > 1$). The healthy resting heartbeat sits near the boundary between the two. On this reading, $\alpha_1 \approx 0.75$ marks the midpoint of the fGN range, and $\alpha_1 \approx 0.5$ marks the loss of short-term correlation. Table~\ref{tab:alpha-landmarks} summarises the landmarks used throughout this thesis.

\begin{table}[ht]
\centering
\caption{Landmark values of the short-term scaling exponent $\alpha_1$.}
\label{tab:alpha-landmarks}
\small
\begin{tabularx}{\textwidth}{l X X}
\toprule
$\alpha_1$ & Signal class & Typical physiological context \\
\midrule
$\approx 1.5$ & Brownian-like, strongly correlated drift & Excessive regularity, e.g.\ slow periodic oscillations dominating the series \citep{goldberger2002} \\
$\approx 1.0$ & $1/f$, boundary of fGN and fBM & Healthy rest and very low exercise intensity \citep{gronwald2020biomarker} \\
$\approx 0.75$ & Midpoint of the fGN range & First ventilatory threshold in incremental exercise \citep{rogers2021threshold} \\
$\approx 0.5$ & Uncorrelated (white noise) & High exercise intensity; loss of organisation \citep{gronwald2020review} \\
\bottomrule
\end{tabularx}
\end{table}

\subsubsection{Window length, detrending and reference settings}

Because \al{} spans only 4 to 16 beats, it can be estimated from short recordings, but the estimate becomes unstable when too few boxes are available. In exercise research, \al{} is commonly computed over a two-minute window, recomputed every few seconds as the window slides along the recording \citep{rogers2021threshold,rogers2022update}. The two-minute convention is supported by analyses of how data length and non-stationarity affect DFA estimates \citep{chen2002} \verify{exact claim in Chen et al.\ on minimum length}. At the heart rates seen in light activity, two minutes correspond to roughly 120 to 200 beats. The present study adopts the two-minute window because it is also close to the longest duration for which a skilled manipulation movement can be sustained without breaking off.

Two further settings affect the estimate, the first of which is the detrending of the RR series before DFA. Reference implementations, including the widely used Kubios software, apply smoothness-priors detrending \citep{tarvainen2002} with a smoothing parameter of $\lambda = 500$ to remove very slow trends that would otherwise inflate the fluctuation function \citep{rogers2022update}. The second is the treatment of artifacts, as missed, extra and ectopic beats distort \al{} more strongly than amplitude indices, and artifact rates above a few percent bias the estimate \citep{rogers2022update}. \citet{hoos2025} argue that these settings, together with window length and box-size range, need to be standardised and reported explicitly before \al{} values can be compared between studies. Section~\ref{sec:analysis} specifies the settings used here.

Taken together, \al{} summarises the short-range correlation structure of the heartbeat in a single exponent with interpretable landmarks, provided that window length, detrending and artifact handling are fixed and reported.

\subsection{Acute Changes in DFA $\alpha_1$}
\label{sec:bg-acute}

\subsubsection{Rest versus activity}

\citet{karasik2002} were among the first to report that correlations in heartbeat fluctuations differ between rest and exercise, with the short-range structure changing more than the long-range structure \verify{direction and scale of the effect in Karasik et al.}. Subsequent exercise research established that \al{} decreases systematically with increasing intensity. A systematic review found this pattern consistently across endurance exercise studies, in which \al{} sat around 1.0 at rest and low intensity, fell through the moderate range and approached 0.5 or below at high intensity \citep{gronwald2020review}. The decline is generally interpreted as a loss of correlation as the multiple regulatory influences on the heart are progressively replaced by a dominant drive from central command and metabolic feedback \citep{gronwald2020biomarker,balague2020}.

\subsubsection{\texorpdfstring{$\alpha_1$}{alpha1} as an intensity marker}

The intensity dependence of \al{} has made it a candidate biomarker for exercise prescription. \citet{rogers2021threshold} proposed that the heart rate at which \al{} crosses 0.75 during an incremental test approximates the first ventilatory threshold, the boundary between low and moderate intensity. Later work broadened this to a general account of \al{} as a marker of intensity distribution in training \citep{gronwald2020biomarker,rogers2022update}. However, the proposal has been challenged. For example, \citet{cassirame2025} examined DFA-derived thresholds in an incremental cycling test and questioned how reliably they match physiological thresholds in individuals \verify{exact findings of Cassirame et al.}. \citet{hoos2025} cautioned that fixed cut-offs depend on reference settings that are not yet standardised. For the present thesis, the relevant point is that \al{} responds within minutes to changes in physiological state.

\subsubsection{Beyond intensity}

A smaller body of work shows that \al{} also responds to factors other than momentary intensity. In ultramarathon runners, \al{} measured at the same low running intensity was lower after the race than before it, suggesting that fatigue alters the organisation of heart rate regulation independently of workload \citep{rogers2021fatigue}. Short-term DFA has also been applied to psychological stress paradigms \todo{find and add a verified source on short-term DFA under cognitive or psychological stress}. These findings matter for the present question because they show that two conditions with the same heart rate can differ in \al. This is the logic on which the heart-rate adjustment in the present design rests.

What is missing from this literature is a task in which physical intensity is low but attentional and motor demands are high. Almost all acute \al{} studies concern endurance exercise, in which intensity and regulatory demand rise together. Object manipulation decouples the two: it raises heart rate only modestly but demands continuous visuomotor prediction and fine force control.

Taken together, \al{} responds within minutes to changes in physiological state, and it can differ between two conditions with the same heart rate. What remains untested is a task that is light on the body and heavy on control.

\subsection{Rhythmic Movement, Sound and Cardiorespiratory Coupling}
\label{sec:bg-coupling}

\subsubsection{Locomotor--respiratory coupling}

Rhythmic movement interacts with the body's other rhythms, and during walking, running, cycling, rowing and wheelchair propulsion, breathing tends to lock into stable frequency ratios with the movement cycle, a phenomenon called locomotor--respiratory coupling (LRC) \citep{bramble1983,bernasconi1993}. In quadrupeds, coupling is enforced mechanically by the loading of the thorax with each stride \citep{bramble1983}; in humans it is weaker and more variable, and its strength differs between individuals and tasks \citep{bernasconi1993}. \citet{bardy2015} distinguish three sources of coupling: mechanical interactions, such as the displacement of the viscera with each step; neurophysiological interactions, such as afferent signals from moving limbs reaching respiratory centres in the brainstem; and informational interactions, such as external rhythms the person attends to.

LRC can be modelled as the interaction of two coupled nonlinear oscillators. \citet{hoffmann2012} used the sine circle map, in which the ratio of the two frequencies (the winding number) and the strength of the coupling determine whether the system settles into a stable mode-locked ratio. The stable ratios are ordered in a Farey tree, in which low-order ratios such as 1:2 occupy wider regions of stability than high-order ratios, and are therefore more resistant to perturbation \citep{hoffmann2012,hoffmann2015}.

\subsubsection{Sound-induced stabilisation}

External sound can strengthen this coupling, as shown in a cycling study in which participants synchronised either their pedalling or their breathing with a metronome set to their own preferred frequency \citep{hoffmann2012}. Auditory pacing stabilised LRC and was accompanied by a significant reduction in energy expenditure. The effect was asymmetrical: pacing the pedalling stabilised coupling more than pacing the breathing. When participants cycled in time with the tone, their breathing shifted towards more stable frequency ratios, as the sine circle map predicts. A follow-up study varied the frequency of pacing \citep{hoffmann2015}. Sound stabilised LRC at all frequencies tested and for both paced systems. The effect was strongest at the participant's preferred rhythm, and frequencies far from it led to a loss of stability. The authors concluded that coupling runs in both directions, from movement to breathing and from breathing to movement.

\citet{bardy2015} reviewed both studies and drew out the role of the starting state. Sound stabilised participants whose coupling was initially unstable, and it could destabilise those whose coupling was already stable \verify{whether this is reported in Bardy et al.\ 2015 or Hoffmann et al.\ 2012}. The effect of rhythmic information therefore depends on how well coordinated the person already is, which is precisely the kind of individual difference that expertise might produce.

\subsubsection{Music, rhythm and the heart}

The cardiovascular effects of music have been studied for decades, with results that range from careful laboratory physiology to broad clinical claims. Narrative reviews conclude that music can modulate heart rate, blood pressure, HRV and stress hormones, but that effects depend on the music and the listener and that the underlying mechanisms remain poorly understood \citep{trappe2010,cervellin2011}. Some of these reviews go further, recommending particular composers and warning against particular genres \citep{trappe2010}; such claims rest on limited evidence and are noted here as context only. A randomised controlled study gives a more measured picture \citep{trappe2016}. Twenty-five minutes of music by Mozart or Strauss lowered blood pressure and heart rate, whereas music by ABBA did not. Serum cortisol decreased in all groups, including a control group resting in silence.

Laboratory work on tempo is more directly relevant. Faster music elicits greater sympathetic arousal, measured as skin conductance, than slower music \citep{dillman2007}. Skin conductance responses are also more frequent during music than during silence, consistent with musical beats capturing attention. \citet{bernardi2009} recorded heart rate, respiration, blood pressure, cerebral blood flow velocity and skin vasomotion while 24 young adults, half of them trained choristers, listened to music with different profiles. Crescendos raised blood pressure and constricted skin blood vessels in proportion to their intensity, and uniform emphasis did the reverse. Musical phrases repeating about every ten seconds, close to the 0.1~Hz rhythm of blood-pressure Mayer waves, entrained cardiovascular variables. Notably, the responses did not differ qualitatively between musicians and non-musicians. A recent review in the Cambridge Elements series integrates such findings into an account of how auditory rhythm interacts with cardiac, respiratory and motor rhythms in perception and action \citep{wright2022}. Wearable sensing now allows these relations to be measured outside the laboratory; for example, an earable system predicted the direction of physiological change during music listening from vital signs alone in about 70\% of cases \citep{butkow2024}.

Rhythmic activity shared between people can couple their hearts as well. When pairs of non-experts vocalised together, their HRV became more coupled during long, synchronous notes than during short notes or at baseline \citep{ruizblais2020}. Most of this coupling was attributable to respiratory sinus arrhythmia, but part of it persisted after the influence of breathing was removed, and the frequency of vocal onsets matched the frequency at which coupling peaked. Coupling did not, however, map clearly onto participants' subjective sense of togetherness.

\subsubsection{Implications for an upper-limb rhythmic task}

Three conclusions follow for the present study. First, rhythmic movement and rhythmic sound both have documented routes to cardiorespiratory rhythms, mainly via breathing and respiratory sinus arrhythmia, with some coupling remaining once breathing is accounted for \citep{ruizblais2020}. Second, the stick's figure-eight motion is a rhythmic upper-limb movement without the impact loading of locomotion, so any coupling it produces is likely to be mainly neurophysiological and informational \citep{bardy2015}. Third, the effect of rhythmic information depends on the starting state of the person, which motivates measuring expertise. The present study records cardiac dynamics only, without a breathing signal or a metronome, and asks whether a downstream signature of coupling is visible there. Coupling serves as one candidate mechanism among several, alongside effort, attention and arousal, and the discussion returns to how it could be isolated in future work.

\subsection{Skill, Expertise and Flow}
\label{sec:bg-skill}

\subsubsection{Object manipulation as a model of skill}

Object manipulation has a place in cognitive neuroscience as a model of skill learning. Three months of juggling practice produced transient increases in grey matter in the mid-temporal visual motion area hMT/V5 and the left posterior intraparietal sulcus, which partly reversed when practice stopped \citep{draganski2004}. Six weeks of juggling training changed white-matter architecture beneath the right posterior intraparietal sulcus \citep{scholz2009}. These studies established juggling as a model of experience-dependent plasticity, although their focus is structural and long-term. They say little about the state of the person while manipulating an object, which is the focus here.

Classical accounts describe skill acquisition in three stages \citep{fitts1967}. In the \emph{cognitive} stage, the learner works out what to do; movement is effortful, error-prone and under explicit attentional control. In the \emph{associative} stage, errors fall and movement becomes more consistent as the learner refines the mapping between intention and outcome. In the \emph{autonomous} stage, movement runs with little conscious control, freeing attention for other demands. A stick-naive participant performing an isolation for the first time is in the cognitive stage; an experienced practitioner performing the same movement is close to autonomy. The same visible movement can therefore correspond to very different internal states, and any physiological measurement of practice has to take expertise into account.

The two movements in this study differ along the same lines. Isolations demand that one point of the stick stay fixed in space while the hand travels around it, which requires continuous visual feedback and fine corrections. Spins are faster and, once learned, more self-sustaining, because the rotation carries momentum, and the hand's role shifts from placing the stick to maintaining a rhythm. Isolations thus load prediction and control, whereas spins load rhythm and timing, and expertise is expected to change both, possibly in different ways.

\subsubsection{The psychophysiology of flow}

Flow denotes a state of complete absorption in an activity whose challenge matches one's skill, accompanied by a merging of action and awareness, a sense of control and a loss of self-consciousness \citep{csikszentmihalyi1990}. Its physiology has mostly been studied with amplitude measures. For example, in expert pianists playing familiar pieces, flow was associated with higher heart rate, deeper breathing and activity of the smiling muscle (zygomaticus major) \citep{demanzano2010} \verify{exact physiological findings in de Manzano et al.}. The authors read this as positive affect combined with effortful attention. In a task matched to participants' skill, the matched condition came with lower HRV and higher salivary cortisol than boredom or overload \citep{keller2011} \verify{direction of HRV and cortisol effects in Keller et al.}. \citet{peifer2014} brought these findings together by showing that flow relates to sympathetic arousal in an inverted U and to parasympathetic activity positively, suggesting co-activation of both autonomic branches during flow.

Neurocognitive accounts point in a compatible direction. Specifically, the transient hypofrontality hypothesis proposes that flow involves a temporary down-regulation of explicit, prefrontally mediated control, allowing implicit, well-practised processes to run with less interference \citep{dietrich2004}. Imaging of experimentally induced flow, using a task adaptively matched to each participant's skill, found reduced activity in the medial prefrontal cortex and amygdala and increased activity in the inferior frontal gyrus and putamen \citep{ulrich2014}. The medial prefrontal cortex and amygdala are precisely the regions implicated in neurovisceral integration (Section~\ref{sec:bg-autonomic}), which gives a plausible anatomical route from a flow-like state to cardiac regulation.

\subsubsection{The case for a complexity marker}

Correlation structure may capture co-activation of competing autonomic inputs, a pattern that amplitude indices summarise poorly. If a state of skilled absorption involves several regulatory loops operating together, its signature should appear in how heart rate variation is organised, which is what \al{} quantifies. Inducing flow is outside the primary aim of the present study, and participants include novices for whom the task is unlikely to produce it. Flow nevertheless motivates the choice of \al{} as the outcome and the inclusion of expertise as a factor. If the organisation of cardiac regulation during manipulation changes with skill, experts are where a flow-related signature would first be expected.

Taken together, skill learning and flow both predict that the same movement can engage cardiac regulation differently depending on expertise. This prediction has not yet been tested with a measure of complexity.

\subsection{Measuring RR Intervals Outside the Laboratory}
\label{sec:bg-h10}

Electrocardiography (ECG) is the reference method for recording RR intervals, but chest-strap sensors are now widely used in research because they are cheap, portable and tolerate movement. The Polar H10 records a single-lead ECG through electrodes in the strap and transmits RR intervals at a resolution of 1/1024~s over Bluetooth Low Energy.

Two validation studies frame its use here. \citet{schaffarczyk2022} compared the H10 with simultaneous ECG in 25 recreational athletes during rest and an incremental cycling test to exhaustion. Agreement was high at rest before and after exercise (e.g.\ $r = 0.95$ and ICC $= 0.95$ for \al{} before exercise) and correlations remained above 0.93 during exercise. During exercise, however, all variables showed a statistically significant bias. For RR intervals and heart rate, bias and limits of agreement were negligible. For \al, they widened with intensity: from a bias of 0.9--8.6\% with limits of agreement of $+11.6$ to $-9.9$\% at low intensity, to limits of $+58.1$ to $-40.9$\% at high intensity. \al{} values in the uncorrelated range, below about 0.5, were the least reliable. \citet{blalock2026} compared the H10 with ECG in 43 healthy young adults during supine rest, paced deep breathing and active standing. Agreement for mean RR interval, RMSSD, SDNN and standard reflex tests was excellent (concordance correlation $\geq 0.99$, mean absolute percentage error below 1\%, 95\% limits of agreement within about $\pm 1.7$~ms). The authors concluded that the H10 is effectively interchangeable with laboratory ECG for these measures.

Two implications follow. For RR intervals and amplitude indices at rest and during light activity, the H10 can be treated as equivalent to ECG. For \al, it is valid at rest and at low intensity, which is the range expected during stick manipulation, but its error grows with intensity and with movement of the upper body, which object manipulation introduces. Artifact rates must therefore be reported per trial, and agreement between the real-time and offline \al{} estimates is itself a question for the present study.

Taken together, the H10 is a valid tool for this study as long as the movement stays light and data quality is checked trial by trial.

\subsection{Research Gap}
\label{sec:bg-gap}

The review identifies three gaps. First, with regard to the \emph{domain}, research on acute changes in \al{} concerns endurance exercise, in which physical intensity and regulatory demand rise together; brief, attention-demanding motor skills at low intensity have not been studied. Second, with regard to the \emph{index}, the psychophysiology of flow and skilled engagement relies on amplitude indices such as RMSSD and high-frequency power, and the organisation of variability during skilled action has not been examined. Third, with regard to \emph{expertise}, both the coupling literature and the skill literature show that the same input or movement affects people differently depending on their starting state. No study has asked whether the cardiac signature of a movement depends on expertise with it.

Underlying all three is a methodological problem. Because \al{} falls with exercise intensity, any change during a movement task could reflect physical load alone. A study of skilled manipulation must therefore treat heart rate as a covariate, so that an effect counts as specific only if it exceeds what heart rate alone predicts. For example, suppose a participant's heart rate rises from 75 to 90 beats per minute during spins and their \al{} drops from 1.05 to 0.85. If the same 15-beat rise during light cycling typically lowers \al{} by about the same amount, the spins add nothing beyond load. If \al{} drops much further than that, something about the skill itself is at work \todo{replace the illustrative numbers with pilot values}.

\subsection{Hypotheses}
\label{sec:bg-hypotheses}

Young, healthy participants are likely to rest near $\alpha_1 \approx 1.0$, which leaves little room for an \emph{increase} to be interpreted as an improvement. I therefore formulate the hypotheses in terms of direction and specificity, and read them with the two-directional complexity framework (Section~\ref{sec:bg-dfa}): (1) a shift towards 0.5 indicates a loss of short-term correlation; (2) a shift towards 1.5 indicates increasing regularity; (3) stability near 1.0 indicates preserved organisation under task demand. All hypotheses were specified before the main sample was collected \confirm{pre-registration with Timo, date and platform}.

Specifically, I hypothesised that mean \al{} during stick manipulation differs from \al{} during a standing baseline in which the participant holds the stick still (Hypothesis~1). As no direction was predicted, this is tested two-sided, and the direction of the shift and the position of task values relative to the landmarks 0.5, 0.75 and 1.0 are reported regardless of significance.

Furthermore, as \al{} falls with exercise intensity, a shift during manipulation may reflect physical load alone. I therefore hypothesised that the shift during manipulation exceeds what the accompanying rise in heart rate predicts, and that this specific component differs between isolations and spins (Hypothesis~2). Spins are faster and more rhythmic, so I expect heart rate to explain most of their shift. Isolations are slow, precise and visually guided, so any shift beyond heart rate would point to regulatory engagement.

Finally, I expected expertise to moderate the shift (Hypothesis~3). Less experienced participants are closer to the cognitive stage of skill, so the same movement demands more explicit control from them. At comparable heart rate, they should therefore depart further from their baseline \al{} than highly practised participants. Because the sample consists mainly of experienced practitioners, expertise is tested as a continuous score, and the few novices are compared with the experienced group descriptively.

I examine three further questions without a directional hypothesis: (1) whether \al{} at the second baseline differs from the first (carry-over); (2) how closely the real-time \al{} of the streaming application agrees with the offline estimate; and (3) whether the dominant and non-dominant hand differ.
